\documentclass[11pt]{article}
\usepackage[margin=1in]{geometry}
\usepackage{amsmath,amssymb}
\usepackage{enumitem}

\title{LEHD Microdata Questions}
\date{}

\begin{document}
\maketitle

\section*{Data Available for Workers}

\paragraph{Earnings (EHF/JHF).}
Quarterly earnings from state UI wage records. Near-universal coverage ({>}95\% of employment; excludes independent contractors, federal employees, railroad workers). About 1.5\% of job-quarters have imputed earnings due to underreporting.

\paragraph{Residence (ICF\_US\_RESIDENCE).}
Annual residential location at census-block level (lat/lon and FIPS geocode).

\paragraph{Sex, date of birth, place of birth (ICF, Stage~A).}
${\sim}5\%$ missing. Imputed via nonparametric KDE using Numident, Census, and ACS as sources.

\paragraph{Race and ethnicity (ICF, Stage~B).}
${\sim}20\%$ missing. Imputed via CART + KDE conditioning on sex, DOB, earnings quartiles, POB, and coworker/co-resident demographics. Moderate quality; the missingness rate is nontrivial but conditioning variables are informative.

\paragraph{Education (ICF, Stage~C).}
${\sim}80$--$85\%$ missing; only available for workers age~25+. Four categories ({$<$}\,HS, HS, some college, BA+). Observed values come from Census 2000 long form (${\sim}7\%$ coverage) and ACS (${\sim}15\%$ cumulative). Imputation uses a log-linear model + KDE with two quality tiers:
\begin{itemize}[nosep]
    \item \emph{Full-information imputation} (workers with pre-2009 earnings): conditions on sex, age, race, POB, earnings deciles, industry, coworker/co-resident demographics. Run once in 2010 and never updated.
    \item \emph{Limited-information imputation} (workers entering 2009+): conditions only on sex, age, race, ethnicity, POB.
\end{itemize}
Define a cell $g$ by the interaction of conditioning variables, with $c \in \{1,2,3,4\}$ indexing education categories. For the full-information tier, cells are sparse, so cell probabilities are estimated via log-linear regression:
\[
    \ln \pi_{g,c} = \mu_c + \sum_k \beta_{k,c}\, x_{g,k} + \text{(selected interactions)}.
\]
The fitted $\hat{\pi}_{g,c}$ parameterize a Dirichlet PPD from which $L=10$ implicates are drawn. For the limited-information tier, cells $g'$ are coarser and use raw counts directly; the coarse-cell probabilities marginalize over the omitted variables:
\[
    \pi_{g',c}^{\text{lim}} = \sum_{g \in g'} \frac{n_g^{\text{obs}}}{n_{g'}^{\text{obs}}}\, \pi_{g,c}^{\text{full}}.
\]
All parameters are frozen from the 2010 run. Out-of-sample accuracy: 26--47\% correct category (vs.\ 80--94\% self-report consistency across Census and ACS).

\paragraph{Employer linkage.}
Workers are linked to establishments (SEIN/SEINUNIT) via UI records. For multi-unit employers, worker-to-establishment assignment is imputed (U2W) using establishment size and worker residence-to-establishment distance.

\section*{Priority Questions for Microdata Side of Model}

\subsection*{1.\ Proxy for worker skill}

Require $(u_{ij})_j \perp v_i \mid (w_j)_j, (\xi_j)_j$ and $\ln q_i$ strictly increasing in $v_i$.

\medskip\noindent\textbf{Past wage:}
\begin{itemize}[nosep]
    \item Transparent, available, and for experienced workers can use previous-firm wage.
    \item Concerns:
    \begin{enumerate}[nosep,label=\alph*.]
        \item \emph{Contamination with location}: my past wage reflects my location and my skill. If $v_i$ is a function of past $(d_{ij})_j$, it is likely that unobserved skill $e_i$ is correlated with current $(d_{ij})_j$; this is an exclusion violation for distance. Concretely, a high-skill worker who lived near low-wage firms may have low past wages to avoid long commutes; this worker's skill residual is large and positive, and correlated with their location.
        \item \emph{Contamination with preference}: my past wage reflects my preference and my skill. If $v_i$ is a function of past $(u_{ij})_j$, it is likely that unobserved skill $e_i$ is correlated with current $(u_{ij})_j$; this is an exclusion violation for $v_i$. Concretely, a high-skill worker who idiosyncratically prefers a low-wage firm may have low past wages; this worker's skill residual is large and positive, and correlated with their preferences.
    \end{enumerate}
    \item Concern~(a) can be somewhat addressed by looking at past wages of my coworkers (my neighbors have past $d_{ij}$ similar to mine), assuming that workers do not systematically sort near firms where they work.
    \item Concern~(b) can be somewhat addressed by looking at past wages of my neighbors (my coworkers have past $\varepsilon_{ij}$ similar to mine), assuming that there is no spatial correlation in $u_{ij}$ (which amounts to assuming that workers do not sort near firms they idiosyncratically prefer).
\end{itemize}
Measures based on auxiliary demographic or income data of my coworkers or neighbors (e.g.\ census tract income measures) share similar properties.

\subsection*{2.\ Within-firm wage variation}

Wages vary within establishment. This poses at least two questions:
\begin{enumerate}[nosep,label=\alph*.]
    \item How many wages and cutoffs does the firm set, and for which groups are these tailored? We could define education $\times$ experience cells and define firm production in terms of an aggregate of different kinds of labor. How do we identify substitution between different kinds of labor?
    \item (For firm level moments) What wage should we use when estimating the wage preference coefficient in $\mathbb{E}[\delta_j - \eta \ln w_j \mid z_j^1] = 0$?
\end{enumerate}

\subsection*{3.\ Stratifying preferences for firms}
Different sorts of workers may have different preferences for firms, say according to education or demographics. This is only an issue for estimation if these preferences are also correlated with skill, i.e.\ we need to add demographic or education subscripts to $\xi_j$ in the exclusion restriction $(u_{ij})_j \perp v_i \mid (w_j)_j, (\xi_j)_j$.

Stratifying preferences by demographics, though, means we cannot use these variables as skill shifters.

%% ---------------------------------------------------------------
%% Commented-out sections (instruments)
%% ---------------------------------------------------------------
% \subsection*{Instrument for firm scale}
% Require $\mathbb{E}[\ln A_j \mid z_j^2] = 0$.
% For example, observed exogenous variation in amenities. Exogenous amenities are an instrument for the relationship between revenue and (skill-composition-adjusted) firm size.
% Ideas:
% \begin{itemize}[nosep]
%     \item Entry or exit of nearby firms
% \end{itemize}

% \subsection*{Instrument for wage}
% Require $\mathbb{E}[\ln \xi_j \mid z_j^1] = 0$.
% For example, observed exogenous variation in TFP\@. Exogenous TFP is an instrument for the relationship between (screening-adjusted) firm size and wages.

\end{document}
